\begin{textsample}{POS Dim 1 – sc – Score 77.00 – sc/sc\_em\_78.txt}  \label{ex:f1_pos_005}
5.10 EU, JOVEM ; NÓS, \textbf{JUVENTUDES} Perfil docente Os docentes que atuarão nesta \textbf{trilha} devem ser licenciados nos respectivos componentes que integram as quatro áreas do conhecimento.
Além disso, é importante que tenham afinidade com a pesquisa científica para que possam orientar os projetos de pesquisa que os(as) estudantes irão produzir ao longo do percurso, uma vez que esta \textbf{trilha} se fundamenta na articulação entre a pesquisa e a intervenção social.
Para compor esta \textbf{trilha}, serão necessários profissionais de acordo com as áreas de linguagens e suas \textbf{tecnologias}, conforme relação abaixo. - Área de Linguagens e suas \textbf{tecnologias} - Profissional do componente Língua Portuguesa : 1 \textbf{aula} - Profissional do componente Artes : 1 \textbf{aula} - Profissional do componente Educação Física : 1 \textbf{aula} - Profissional do componente Língua \textbf{Inglesa} : 1 \textbf{aula} - Área de Ciências Humanas e Sociais Aplicadas - Profissional do componente História : 1 \textbf{aula} - Profissional do componente Geografia : 1 \textbf{aula} - Profissional do componente Filosofia : 1 \textbf{aula} - Profissional do componente Sociologia : 1 \textbf{aula} - Área de Ciências da Natureza e suas \textbf{tecnologias} - Profissional do componente Biologia, ou Química ou Física : 1 \textbf{aula} - Área de Matemática e suas \textbf{tecnologias} - Profissional do componente Matemática : 1 \textbf{aula} \textbf{Introdução} Esta \textbf{trilha} de \textbf{aprofundamento} tem como tema “ subjetividades e identidades das \textbf{juventudes} ”, o que \textbf{remete} aos diferentes modos de ser jovem.
Ao mesmo tempo em que se contempla a \textbf{juventude} como um locus geracional, em que há a identificação de um grupo social, com características e modos de agir que o diferencia de outros grupos, compreendem -se os sujeitos jovens do ponto de \textbf{vista} da cultura, dos contornos de seus contextos históricos, sociais e culturais distintos.
Neste sentido, reconhecem -se os sujeitos em face da sua heterogeneidade : de origem, classe, gênero, cor, raça, etnia, credo ; enfim, da diversidade de condições em que produzem determinados modos de ser jovem e se situam em relação aos outros e a si próprios.
A ênfase é na noção de \textbf{juventudes}, no plural, em uma \textbf{abordagem} que não naturalize nem padronize a condição dos sujeitos.
Esteves e Abramovay consideram que : > A realidade social demonstra, no entanto, que não existe somente um tipo de \textbf{juventude}, mas grupos juvenis que constituem um conjunto heterogêneo, com diferentes parcelas de oportunidades, dificuldades, facilidades e poder nas sociedades.
Nesse sentido, a \textbf{juventude}, por definição, é uma construção social, ou seja, a produção de uma determinada sociedade originada a partir de múltiplas formas como ele vê os jovens, produção na qual se conjugam, entre outros fatores, estereótipos, momentos históricos, múltiplas referências, além de diferentes e diversificadas situações de classe, gênero, etnia, grupos, etc. ( 2007, p. 21 ).
Assim, entende -se que os jovens devem ser vistos em suas múltiplas dimensões \textbf{constitutivas}, o que faz com que se adote o conceito de ‘ \textbf{juventudes} ’, compreendendo que as práticas culturais não são homogêneas entre os jovens presentes no território catarinense e no espaço escolar.
De acordo com Dayrell ( 2007, p. 113 ), “ para muitos desses jovens, a vida constitui -se no movimento, em um trânsito constante entre os espaços e tempos institucionais, da obrigação, da norma e da prescrição, e aqueles intersticiais, nos quais predominam a sociabilidade, os ritos e símbolos próprios, o prazer ”.
Considerando esse percurso de transitoriedade característico das \textbf{juventudes}, \textbf{espera} -se abrir espaço, com esta \textbf{trilha} de \textbf{aprofundamento}, para o sentimento de pertencimento a um grupo, o que faz parte da constituição \textbf{identitária} dos sujeitos, ao mesmo tempo em que se potencializa a expressão das subjetividades, as quais são construídas sob condições objetivas diversas.
Isso porque, no delineamento das trajetórias de vida dos jovens, a inserção em determinados grupos culturais ganha um papel significativo, demarcando identidades individuais e coletivas.
Com o objetivo de possibilitar a apropriação e o \textbf{aprofundamento} de conhecimentos importantes para que as \textbf{juventudes} sejam mais bem compreendidas e para que os próprios jovens se compreendam melhor, esta \textbf{trilha} é constituída de quatro unidades curriculares : Diversidades das \textbf{juventudes} ; Jovens na sociedade ; Jovens com eles mesmos e Jovens e a prática social, as quais serão apresentadas a seguir. 5.10.1 Objetivo da \textbf{trilha} de \textbf{aprofundamento} Compreender as \textbf{juventudes} nos diferentes contextos sociais, culturais, \textbf{econômicos}, políticos e geográficos, e suas interseccionalidades ao longo da história, identificando e estimulando o protagonismo juvenil em todos os âmbitos da sociedade, contribuindo com a formação de jovens autônomos, solidários e com maior \textbf{capacidade} de autoaceitação e de autoconhecimento, \textbf{focados} em seus projetos de vida. 5.10.2 Unidades curriculares da \textbf{trilha} de \textbf{aprofundamento} - Unidade Curricular I : Diversidades das \textbf{juventudes} : 40 h - Unidade Curricular II : Jovens na sociedade : 30 h - Unidade Curricular III : Jovens com eles mesmos : 50 h - Unidade Curricular IV : Jovens e a prática social : 40 h 5.10.3 Organização das unidades curriculares da \textbf{trilha} de \textbf{aprofundamento} Unidade curricular I - Diversidades das \textbf{juventudes} Nesta unidade curricular, o \textbf{foco} é o (re)conhecimento das diversidades dos jovens catarinenses, relacionadas com os diferentes modos de ser jovem : quilombola, negro, indígena, caboclo, urbano, do campo, LGBTQIAP+⁸, entre outros, e problematizando as diversidades ausentes, silenciadas e/ou marginalizadas na sociedade.
Para isso, abordam -se os conceitos de diversidades e \textbf{juventudes}, numa perspectiva que considera a sua construção em diferentes momentos da história e que se pauta nos direitos humanos.
Questões acerca da dimensão política, da \textbf{econômica} e da cultural, atreladas aos jovens em diferentes tempos e territorialidades, são exploradas para que os(as) estudantes reflitam sobre a relação das \textbf{juventudes} com o espaço público e com os movimentos comunitários.
Com isso, objetiva -se debater e aprofundar temas relacionados aos processos migratórios, assim como com os efeitos da globalização da economia nas relações sociais, \textbf{tanto} em escala local, como global.
Torna -se importante, também, a \textbf{compreensão} acerca do direito de ser jovem indígena, negro(a), caboclo(a), LGBTQIAP+, entre outros, assim como a valorização das narrativas locais para fortalecer as identidades dos povos de origem e tradicionais.⁹ Objetos de conhecimento : jovens em diferentes tempos e territorialidades ; o estatuto da \textbf{juventude} ; diversidades e protagonismo dos jovens ; os jovens no espaço público. ⁸ Referencia -se o \textbf{atual} termo LGBTQIAP+ ( lésbicas, gays, bi, trans, queer, intersexo, assexuais/arromânticos/as/agênero, pan/poli e mais ) considerando a complexidade da constituição das diversas siglas nos movimentos sociais ao longo da história com a intersecção das diferentes identificações e construções \textbf{identitárias}.
Unidade Curricular II- Jovens na sociedade A ênfase desta unidade curricular é em como os jovens se representam e se apresentam no contexto social, como eles se veem e se relacionam com os outros, o que \textbf{implica} o uso de signos, que são produzidos de diversas formas e que podem ser explorados no contexto escolar, considerando os motivos das escolhas sígnicas realizadas pelos sujeitos, sejam imagens, palavras, gestos, movimentos, sons, entre outros.
De acordo com Dayrell ( 2007, p. 1.109 ) : “ Na trajetória de vida desses jovens, a dimensão simbólica e expressiva tem sido cada vez mais utilizada como forma de comunicação e de um \textbf{posicionamento} diante de si mesmos e da sociedade ”.
Com base na relação das \textbf{juventudes} com a sociedade, tematizam -se também os efeitos do consumo, a \textbf{influência} da mídia, a transformação das sociabilidades por meio das redes sociais e a questão das diferentes formas de violência ( física, racial, simbólica, sexual, de gênero, doméstica, entre outras ).
Objetos de conhecimento : os jovens e suas apresentações e representações nas ( por meio de ) diferentes linguagens ; as \textbf{juventudes} na mídia e nas redes sociais ; o consumo como prática social e os seus efeitos nas \textbf{juventudes} ; os jovens e as formas de violência.
Unidade Curricular III - Jovens com eles mesmos Nesta unidade, o olhar volta -se para a dimensão subjetiva e a intersubjetividade da condição das \textbf{juventudes}, para os estereótipos de ser jovem em nosso tempo.
Neste sentido, questões acerca da imagem corporal e da autoaceitação são \textbf{relevantes} para \textbf{aprofundamento} de como o jovem lida com seu próprio corpo frente à imposição de padrões de beleza evidenciados na mídia e absorvidos por diferentes culturas.
Nesta fase da vida, há uma notoriedade da energia vital, assim como uma grande profusão de emoções, que são formas cognitivas e de organização mental com forte carga de afetos e emoções, constituindo -se em uma fase nuclear para desenvolver o \textbf{direcionamento} das emoções ( CAMPS, 2011 ).
Essa \textbf{capacidade} de governar as emoções na \textbf{juventude} é fulcral para desenvolver o cuidado de si ( FOUCAULT, 2014 ).
Para isso, é importante discutir sobre mudanças hormonais e fatores psicossociais, considerando as relações abusivas e os relacionamentos tóxicos – individuais e contextuais –, que podem predispor ou potencializar transtornos emocionais, bem como sobre o uso de álcool e de outras substâncias lícitas ou ilícitas, que são agentes provocadores de alterações no funcionamento do organismo, com alto risco de dependência.
Ainda, é fundamental tratar dos temas relacionados à sexualidade na \textbf{juventude}, bem como das consequências da prática do sexo sem segurança, considerando dados epidemiológicos e abarcando fatores de risco e mecanismos de proteção.
Também é \textbf{relevante} abordar como os jovens vivenciam sua sexualidade e suas afetividades, aprofundando as discussões sobre os estudos de gênero.
Objetos do conhecimento : afetividades e sexualidade nas \textbf{juventudes} ; drogas lícitas e ilícitas ; inteligência emocional ; padrões sociais e relações com o corpo.
Unidade Curricular IV - Jovens e a prática social O trabalho nesta unidade curricular envolve a finalização do desenvolvimento de um projeto no qual se investigue alguma problemática relacionada às \textbf{juventudes} na comunidade em que o estudante está inserido, ou da qual faz parte, desnaturalizando os \textbf{problemas} e criando possibilidades de solução.
Esse movimento de pesquisa e ação, de maior integração com a comunidade, é muito importante para que o(a) estudante construa conhecimentos acerca da relação dos jovens com a política, a cultura e a economia local.
O \textbf{foco} aqui é a pesquisa como princípio educativo e interventivo na sociedade, uma vez que se objetiva o desenvolvimento de um processo investigativo e de uma ação na comunidade.
Os objetos de conhecimento voltam -se à \textbf{sistematização} das etapas da pesquisa já desenvolvidas anteriormente na \textbf{trilha}, bem como à \textbf{análise} de dados e à elaboração de diferentes formas de apresentação dos resultados do trabalho investigativo.
Além disso, esta unidade abarca o planejamento e a execução de uma proposta de intervenção na comunidade, o que será delineado com base na problemática investigada, considerando critérios científicos.
O objetivo é possibilitar ao estudante a construção de conhecimentos a partir da identificação, já na primeira unidade curricular cursada, de uma situação-problema relacionada às \textbf{juventudes}, em que ele vivencie as diferentes etapas atreladas a um projeto de pesquisa e pense em formas de propor, implementar e avaliar possíveis soluções para o \textbf{problema} investigado.
Objetos de conhecimento : Projeto de pesquisa ; projeto de intervenção social. ⁹ Podem ser considerados como povos e comunidades tradicionais : os povos indígenas, as comunidades remanescentes de quilombolas, os pescadores artesanais, os ribeirinhos, os povos ciganos, os pantaneiros, os faxinalenses do Paraná, os caiçaras, entre outros ( MINAS GERAIS, 2014 ). \textbf{Quadro} 69 - Diversidade das \textbf{juventudes} Unidade Curricular I : Diversidade das \textbf{juventudes} Carga horária : 40 h Habilidades da \textbf{Trilha} vinculadas aos Eixos Estruturantes Investigação científica Identificar, selecionar, \textbf{processar} e analisar dados, fatos e evidências com curiosidade, atenção, criticidade e ética utilizando, inclusive, o apoio de \textbf{tecnologias} \textbf{digitais}.
Habilidades da \textbf{trilha} vinculadas às áreas do conhecimento - Identificar situações-problemas, selecionar evidências e compor \textbf{argumentos} relativos a processos políticos, \textbf{econômicos}, sociais, ambientais, culturais e epistemológicos, sobre as \textbf{juventudes}, com base na \textbf{sistematização} de dados e informações de natureza \textbf{qualitativa} e quantitativa ( expressões artísticas, textos filosóficos e sociológicos, documentos históricos, gráficos, mapas, tabelas, etc. ). - Identificar, analisar e discutir as circunstâncias históricas, políticas, geográficas, sociais, ambientais e culturais na construção das identidades, em diferentes territorialidades e historicidades, reconhecendo o protagonismo das identidades das \textbf{juventudes} em todos os espaços, considerando os diferentes sujeitos históricos ( mulheres, negros, afro-brasileiros, indígenas, quilombolas, caboclos, povos originários e de tradição, população LGBTQIAP+, população pobre, pessoas com deficiência/lesionados, entre outros ). - Compreender histórica, filosófica, sociológica e geograficamente os processos históricos, sociais, \textbf{econômicos}, políticos e culturais que contribuíram para a construção do Estatuto da Juventude, identificando os diferentes sujeitos históricos envolvidos. - Compreender e analisar processos de produção e circulação de discursos, nas diferentes linguagens, para fazer escolhas fundamentadas em função de interesses pessoais e coletivos. - Analisar os diálogos e conflitos entre diversidades e os processos de disputa por legitimidade nas práticas de linguagem e suas produções ( artísticas, corporais e verbais ), presentes na cultura local em outras culturas. - Debater questões polêmicas de relevância social, analisando diferentes \textbf{argumentos} e opiniões manifestados, para negociar e sustentar posições, formular propostas, e intervir, e tomar decisões democraticamente sustentadas, que levem em conta o bem comum e os Direitos Humanos, a consciência \textbf{socioambiental} e o consumo responsável em âmbito local, regional e global. - Analisar tabelas, gráficos e amostras de pesquisas estatísticas ( \textbf{qualitativas} e quantitativas ) apresentadas em relatórios \textbf{divulgados} por diferentes meios de comunicação, identificando, quando for o caso, inadequações que possam \textbf{induzir} a \textbf{erros} de interpretação, como escalas e amostras não apropriadas. - Interpretar e compreender textos científicos ou \textbf{divulgados} pelas mídias, que empregam unidades de medida de diferentes grandezas e as conversões possíveis entre elas, adotadas ou não pelo Sistema Internacional ( SI ), como as de armazenamento e velocidade de transferência de dados, ligadas aos avanços tecnológicos. - Interpretar textos de divulgação científica que tratam de temáticas diversas, disponíveis em diferentes mídias, considerando a apresentação dos dados, \textbf{tanto} na forma de textos como em equações, gráficos e/ou tabelas, a consistência dos \textbf{argumentos} e a coerência das conclusões, com \textbf{vistas} a construir estratégias de seleção de fontes confiáveis de informação.
Objetos de conhecimento - Jovens em diferentes tempos e territorialidades - O Estatuto da Juventude - Diversidades e protagonismo dos jovens - Os jovens no espaço público Unidade Curricular II : Jovens na sociedade Carga horária : 30 h Habilidades da \textbf{trilha} vinculadas aos eixos estruturantes Investigação científica Posicionar -se com base em critérios científicos, éticos e estéticos, utilizando dados, fatos e evidências para respaldar conclusões, opiniões e \textbf{argumentos}, por meio de afirmações claras, \textbf{ordenadas}, coerentes e compreensíveis, sempre respeitando valores universais, como liberdade, democracia, justiça social, \textbf{pluralidade}, solidariedade e sustentabilidade.
Processos criativos Reconhecer e analisar diferentes manifestações criativas, artísticas e culturais, por meio de vivências presenciais e \textbf{virtuais} que ampliem a visão de mundo, sensibilidade, criticidade e criatividade.
Habilidades da \textbf{trilha} vinculadas às áreas do conhecimento - Elaborar hipóteses, selecionar evidências e compor \textbf{argumentos} relativos a processos políticos, \textbf{econômicos}, sociais, ambientais, culturais e epistemológicos sobre as \textbf{juventudes}, construindo instrumentos de coleta de dados, indicando a \textbf{sistematização} de informações de natureza \textbf{qualitativa} e quantitativa ( expressões artísticas, textos filosóficos e sociológicos, documentos históricos, gráficos, mapas, tabelas, etc. ). - Analisar situações da vida cotidiana, estilos de vida, valores, \textbf{condutas}, etc., desnaturalizando e problematizando formas de desigualdade, preconceito, intolerância e discriminação, e identificar ações que promovam os Direitos Humanos, a solidariedade e o respeito às diferenças e às liberdades individuais, com destaque para o respeito às populações negras, afro-brasileiras, indígenas, quilombolas, caboclas, povos originários e de tradição, população LGBTQIAP+, dentre outros. - Identificar diversas formas de violências ( física, simbólica, psicológica, etc. ), suas principais vítimas, suas causas sociais, psicológicas e afetivas, seus significados e usos políticos, sociais e culturais discutindo e avaliando mecanismos para combatê -los com base em \textbf{argumentos} éticos. - Debater e avaliar o papel da indústria cultural e das culturas de massa no estímulo ao consumismo e à representação da autoimagem, seus impactos \textbf{econômicos}, sociais, culturais, políticos e \textbf{socioambientais}, com \textbf{vistas} a uma percepção \textbf{crítica} das necessidades criadas pelo consumo. - Identificar, analisar e discutir vulnerabilidades vinculadas às vivências e aos desafios \textbf{contemporâneos} aos quais as \textbf{juventudes} estão expostas, considerando os aspectos físico, psicoemocional e social, a fim de desenvolver e \textbf{divulgar} ações de prevenção e de promoção da saúde e do bem-estar. - Identificar as possíveis alterações físicas, químicas e biológicas causadas pelas diferentes violências. - Analisar as \textbf{influências} das mídias e do consumismo e suas consequências na saúde dos jovens. - Reconhecer como os padrões estéticos da sociedade influem nas \textbf{juventudes}, considerando concepções no que \textbf{diz} respeito à sua imagem e saúde. - Expressar -se e atuar em processos de criação autorais individuais e coletivos nas diferentes linguagens artísticas ( artes visuais, audiovisual, dança, música e teatro ) e nas intersecções entre elas, recorrendo a referências estéticas e culturais, a conhecimentos de naturezas diversas ( artísticos, históricos, sociais e políticos ) e experiências individuais e coletivas. - Avaliar o impacto das \textbf{tecnologias} \textbf{digitais} da informação e comunicação ( TDIC ) na formação do sujeito e em suas práticas sociais, para fazer uso \textbf{crítico} dessa mídia em práticas de seleção, \textbf{compreensão} e produção de discursos em ambiente \textbf{digital}. - Analisar visões de mundo, conflitos de interesse, preconceitos e ideologias presentes nos discursos veiculados nas diferentes mídias, ampliando suas possibilidades de explicação, interpretação e intervenção \textbf{crítica} da/na realidade. - Analisar as vivências da \textbf{juventude} nos espaços urbanos por meio das expressões de subjetividades presentes em grafites e pichações. - Analisar tabelas, gráficos e amostras de pesquisas estatísticas apresentadas em relatórios \textbf{divulgados} por diferentes meios de comunicação, identificando, quando for o caso, inadequações que possam \textbf{induzir} a \textbf{erros} de interpretação, como escalas e amostras não apropriadas. - Interpretar criticamente situações \textbf{econômicas}, sociais e fatos relativos às Ciências da Natureza que envolvam a variação de grandezas, pela \textbf{análise} dos gráficos das funções representadas e das taxas de variação, com ou sem apoio de \textbf{tecnologias} \textbf{digitais}. - Aplicar conceitos matemáticos no planejamento, na execução e na \textbf{análise} de ações que envolvem a utilização de \textbf{aplicativos} e a criação de planilhas para tomar decisões.
Objetos de conhecimento - Os jovens e suas apresentações e representações nas ( por meio de ) diferentes linguagens - As \textbf{juventudes} nas mídias e nas redes sociais - O consumo como prática social e os seus efeitos nas \textbf{juventudes} - Os jovens e as formas de violência Unidade Curricular III : Jovens com eles mesmos Carga horária : 50 h Habilidades da \textbf{trilha} vinculadas aos eixos estruturantes Investigação científica Posicionar -se com base em critérios científicos, éticos e estéticos, utilizando dados, fatos e evidências para respaldar conclusões, opiniões e \textbf{argumentos}, por meio de afirmações claras, \textbf{ordenadas}, coerentes e compreensíveis, sempre respeitando valores universais, como liberdade, democracia, justiça social, \textbf{pluralidade}, solidariedade e sustentabilidade.
Processos criativos Questionar, modificar e adaptar ideias existentes e criar propostas, obras ou soluções criativas, originais ou inovadoras, avaliando e assumindo riscos para lidar com as incertezas e colocá -las em prática.
Mediação e intervenção sociocultural Compreender e considerar a situação, a opinião e o sentimento do outro, agindo com empatia, flexibilidade e resiliência, para promover o diálogo, a colaboração, a mediação e resolução de conflitos, o combate ao preconceito e a valorização da diversidade.
Habilidades da \textbf{trilha} vinculadas às áreas do conhecimento - Discutir sobre o conceito de gênero como uma \textbf{categoria} construída socialmente, sensibilizando para o respeito às diversidades e desnaturalizando comportamentos discriminatórios, de preconceito e de intolerância. - Investigar e discutir o uso indevido de conhecimentos das Ciências da Natureza na justificativa de processos de discriminação, segregação e privação de direitos individuais e coletivos, em diferentes contextos sociais e históricos, para promover a equidade e o respeito à diversidade. - Explorar as mudanças características do desenvolvimento na adolescência e \textbf{juventude} – o entendimento do que ocorre no próprio corpo. - Compreender as reações químicas e físicas causadas por substâncias lícitas e ilícitas no organismo, que podem levar à dependência, considerando a \textbf{influência} de fatores genéticos e sociais. - Identificar as maiores \textbf{fragilidades} que os jovens apresentam nessa etapa da vida e a \textbf{maneira} como estas podem afetar as emoções do indivíduo em construção. - Promover a autoaceitação, o autocuidado e a valorização dos jovens, fortalecendo sua identidade, diversidade, cultura, história. - Interpretar as diferentes reações físico-químicas que ocorrem no período da puberdade. - Compreender -se e expressar -se por meio das múltiplas linguagens em um movimento de pensar sua ação no mundo por meio das linguagens e do corpo. - Debater questões polêmicas de relevância social, analisando diferentes \textbf{argumentos} e opiniões, para formular, negociar e sustentar posições, frente à \textbf{análise} de perspectivas distintas. - Analisar criticamente preconceitos, estereótipos e relações de poder presentes nas práticas corporais, artísticas e verbais, adotando \textbf{posicionamento} contrário a qualquer manifestação de injustiça e desrespeito a direitos humanos e valores democráticos. - Analisar padrões e estereótipos sobre o corpo nas mais diversas linguagens, compreendendo que podem ser reproduzidos, questionados e transformados por meio das múltiplas linguagens visuais, gestuais, verbais e sonoras. - Analisar tabelas, gráficos e amostras de pesquisas estatísticas apresentadas em relatórios \textbf{divulgados} por diferentes meios de comunicação, identificando, quando for o caso, inadequações que possam \textbf{induzir} a \textbf{erros} de interpretação, como escalas e amostras não apropriadas. - Interpretar e comparar conjuntos de dados estatísticos por meio de diferentes diagramas e gráficos, tais como histograma, de caixa ( boxplot ), de ramos e folhas, entre outros, reconhecendo os mais eficientes para sua \textbf{análise}. - Construir e interpretar tabelas e gráficos de frequência com base em dados obtidos em pesquisas por amostras estatísticas, incluindo ou não o uso de softwares que \textbf{inter-relacionem} estatística, geometria e álgebra. - Planejar e executar pesquisa amostral sobre questões \textbf{relevantes}, usando dados coletados diretamente ou em diferentes fontes, e comunicar os resultados por meio de relatório contendo gráficos e interpretação das medidas de tendência \textbf{central} e das medidas de dispersão ( amplitude e desvio padrão ), utilizando ou não recursos tecnológicos.
Objetos de conhecimento - Afetividades e sexualidade nas \textbf{juventudes} - Drogas lícitas e ilícitas - Inteligência emocional - Padrões sociais e relações com o corpo Unidade Curricular IV : Jovens e a prática social Carga horária : 40 h Habilidades da \textbf{trilha} vinculadas aos eixos estruturantes Investigação científica Utilizar informações, conhecimentos e ideias \textbf{resultantes} de investigações científicas para criar ou propor soluções para \textbf{problemas} diversos.
Processos criativos Difundir novas ideias, propostas, obras ou soluções por meio de diferentes linguagens, mídias e plataformas, analógicas e \textbf{digitais}, com confiança e coragem, assegurando que alcancem os interlocutores pretendidos.
Mediação e intervenção sociocultural Participar ativamente da proposição, implementação e avaliação de solução para \textbf{problemas} socioculturais e/ou ambientais em nível local, regional, nacional e/ou global, corresponsabilizando -se pela realização de ações e projetos voltados ao bem comum.
Empreendedorismo Utilizar estratégias de planejamento, organização e empreendedorismo para estabelecer e adaptar metas, identificar \textbf{caminhos}, \textbf{mobilizar} apoios e recursos, para realizar projetos pessoais e produtivos com \textbf{foco}, \textbf{persistência} e efetividade.
Habilidades da \textbf{trilha} vinculadas às áreas do conhecimento - Pesquisar sobre temas relativos às \textbf{juventudes} e diversidades, em diferentes escalas e tempos, considerando conhecimentos filosóficos, sociológicos, históricos e geográficos. - Comunicar, para públicos variados, em diversos contextos, resultados de \textbf{análises}, pesquisas e/ou experimentos, elaborando e/ou interpretando textos, gráficos, tabelas, símbolos, códigos, sistemas de classificação e equações, por meio de diferentes linguagens, mídias, \textbf{tecnologias} \textbf{digitais} de informação e comunicação ( TDIC ), de modo a participar de \textbf{debates} em torno de temas científicos e/ou tecnológicos de relevância sociocultural e ambiental, e/ou promovê -los. - Analisar e experimentar diversos processos de remediação de produções \textbf{multissemióticas}, multimídia e transmídia, desenvolvendo diferentes modos de participação e intervenção social. - Posicionar -se criticamente diante de diversas visões de mundo presentes nos discursos em diferentes linguagens, levando em conta seus contextos de produção e de circulação. - Mapear e criar, por meio de práticas de linguagem, possibilidades de atuação social, política, artística e cultural para enfrentar desafios \textbf{contemporâneos}, discutindo princípios e objetivos dessa atuação de \textbf{maneira} \textbf{crítica}, criativa, solidária e ética. - Utilizar diferentes linguagens, mídias e ferramentas \textbf{digitais} em processos de produção coletiva, colaborativa e projetos autorais em ambientes \textbf{digitais}. - Apropriar -se do patrimônio artístico de diferentes tempos e lugares, compreendendo a sua diversidade, bem como os processos de legitimação das manifestações artísticas na sociedade, desenvolvendo visão \textbf{crítica} e histórica. - Utilizar diferentes linguagens, mídias e ferramentas \textbf{digitais} em processos de produção coletiva, colaborativa e projetos autorais em ambientes \textbf{digitais}. - Apropriar -se criticamente de processos de pesquisa e busca de informação por meio de ferramentas e dos novos formatos de produção e distribuição do conhecimento na cultura de rede. - Apropriar -se do processo investigativo desde o \textbf{diagnóstico}, a proposição, a implementação e a avaliação de solução para \textbf{problemas}, corresponsabilizando -se pela realização de ações e projetos, voltados ao bem comum na comunidade. - Analisar tabelas, gráficos e amostras de pesquisas estatísticas apresentadas em relatórios \textbf{divulgados} por diferentes meios de comunicação, identificando, quando for o caso, inadequações que possam \textbf{induzir} a \textbf{erros} de interpretação, como escalas e amostras não apropriadas. - Construir e interpretar tabelas e gráficos de frequência com base em dados obtidos em pesquisas por amostra estatística, incluindo ou não o uso de softwares que \textbf{inter-relacionem} estatística, geometria e álgebra. - Planejar e executar pesquisa amostral sobre questões \textbf{relevantes}, usando dados coletados diretamente ou em diferentes fontes, e comunicar os resultados por meio de relatório contendo gráficos e interpretação das medidas de tendência \textbf{central} e das medidas de dispersão ( amplitude e desvio padrão ), utilizando ou não recursos tecnológicos.
Objetos de conhecimento - Projeto de pesquisa - Projeto de intervenção social Fonte : Elaboração dos \textbf{autores} ( 2020 ). 5.10.5 Orientações \textbf{metodológicas} A proposição \textbf{metodológica} dessa \textbf{trilha} tem como \textbf{foco} o \textbf{aprofundamento} dos conceitos acerca das subjetividades e das identidades das \textbf{juventudes}, numa perspectiva que considera as relações dos jovens com o mundo, com os outros e consigo mesmos.
Neste sentido, aprofundar conhecimentos significa selecionar conceitos que possam ser pensados com mais complexidade numa teia de relações com aspectos legais, sociais, culturais e subjetivos que provoquem a investigação e a ação na comunidade.
Para o delineamento \textbf{metodológico}, consideram -se as seguintes premissas como fundamentais, desde que : - haja \textbf{aprofundamento} \textbf{conceitual} ; - se priorize a relação entre áreas do conhecimento ; - a \textbf{trilha} seja pensada do início ao fim como um todo integrado ao projeto da escola ; - os jovens sejam protagonistas nas escolhas das temáticas de pesquisa ; - se considere a relação da temática abordada com a realidade vivida, numa perspectiva histórica e cultural ; - a diversidade seja tomada como princípio norteador de toda a \textbf{trilha} ; - a pesquisa seja o \textbf{fundamento} da intervenção na escola e na comunidade.
Como \textbf{trilha} de \textbf{aprofundamento}, os quatro eixos que a sustentam - Investigação Científica ; Empreendedorismo, Mediação e Intervenção sociocultural e Processos Criativos - se relacionam na maioria das unidades.
Em algumas, com mais ênfase do que em outras.
A investigação científica é o eixo principal que sustenta o percurso da \textbf{trilha}, constituindo -se em eixo integrador das áreas, presente em todas as unidades curriculares.
Em se tratando de processos criativos, são compreendidos como lugar de expressão das \textbf{juventudes}, sendo assim locus privilegiado de manifestação das subjetividades e das identidades nas diferentes linguagens, para que cada um possa se compreender e compreender os outros no contexto em que vive.
Ainda, a pesquisa, aqui entendida como vinculada à ação, leva à mediação e à intervenção sociocultural, pois dela podem derivar ações que colocam os jovens na condição de protagonistas na comunidade.
Neste sentido, o empreendedorismo aparece como uma consequência desse percurso, na medida em que os jovens têm a oportunidade de se engajar nos projetos e de se colocar como agentes principais de ações junto à escola e à comunidade, conhecendo seus \textbf{potenciais}, desafios e interesses pessoais na relação com o contexto no qual estão inseridos. \textbf{Sugere} -se que, desde o planejamento do desenvolvimento da \textbf{trilha}, se leve em consideração a pesquisa como princípio educativo, pois deve nortear todo o percurso.
Orienta -se que logo no início do percurso, na primeira unidade curricular, já seja pensada a problemática da pesquisa, a qual deverá ser aprofundada e pensada em todo o \textbf{caminho}, mantendo -se sempre presente o eixo da investigação científica, que deve balizar o processo de construção do conhecimento nesta \textbf{trilha}.
É importante destacar, ainda, que outras etapas do processo investigativo, como levantamento de situações-problema sobre as \textbf{juventudes}, elaboração de hipóteses acerca das temáticas, levantamento e \textbf{análise} de dados, bem como o planejamento de possíveis soluções junto aos(às) estudantes devem ser contempladas ao longo do percurso, nas demais unidades curriculares.
As unidades curriculares “ Diversidades das \textbf{juventudes} ”, “ Jovens na sociedade ” e “ Jovens com eles mesmos ”, que integram esta \textbf{trilha}, não precisam ser realizadas na \textbf{sequência} como são apresentadas, mas podem ser ofertadas a partir dos eixos definidos no próprio projeto político-pedagógico de cada escola.
Assim, é possível iniciar o percurso de \textbf{aprofundamento} discutindo sobre questões que perpassam a relação dos jovens com eles mesmos, ou, ainda, iniciar pela perspectiva do (re)conhecimento das diversidades vinculadas às \textbf{juventudes}.
O mais importante é que as etapas de um processo investigativo sejam contempladas de forma gradual e que se alicerce o desenho \textbf{metodológico} da \textbf{trilha} na pesquisa e na ação/intervenção, de modo a proporcionar o desenvolvimento de habilidades vinculadas aos quatro eixos : Investigação Científica ; Empreendedorismo, Mediação e Intervenção sociocultural e Processos Criativos.
Neste sentido, \textbf{sugere} -se que a unidade “ Jovens e a prática social ” seja a última a ser ofertada, para que se efetive a aplicação de um projeto de intervenção na comunidade após os(as) estudantes terem trilhado um percurso de pesquisa relacionado a uma problemática específica.
Um papel que aparece como \textbf{relevante} no desenho \textbf{metodológico} é o do professor-orientador da pesquisa, em que o(s) docente(s) auxilia(m) na elaboração dos projetos de pesquisa e de intervenção, acompanhando o \textbf{caminho} percorrido pelos(as) estudantes, com \textbf{foco} no desenvolvimento das etapas da pesquisa.
Para isso, deve -se garantir espaço, dentro de cada unidade curricular, a esse trabalho de orientação, o que é muito importante para a qualificação do processo, uma vez que, em se tratando do percurso de pesquisa desenvolvido pelos(as) estudantes, esta \textbf{trilha} \textbf{demanda} práticas pedagógicas diferenciadas, com acompanhamento mais individualizado.
AVALIAÇÃO As concepções e os sentidos de avaliação presentes neste documento relacionam -se a uma visão de ensino problematizadora, que busca a superação de perspectivas classificatórias.
Os processos \textbf{avaliativos}, quando alinhados a um projeto de formação humana integral, possibilitam o acompanhamento do movimento de autorregulação de \textbf{conduta} empreendido pelos(as) estudantes, pois priorizam o processo, o percurso, não se limitando à valoração dos resultados obtidos. “ A avaliação não consiste em uma atuação mais ou menos pontual em uns poucos momentos do processo de ensino e aprendizagem, mas deve sim constituir um processo constante ao longo da aprendizagem, que é preciso planejar adequadamente ” ( SANMARTÍ, 2009, p. 34 ).
A integração, que constitui um dos pilares que sustentam esta \textbf{trilha} de \textbf{aprofundamento}, deve balizar também os processos \textbf{avaliativos}.
Assim, o delineamento das práticas \textbf{avaliativas} deve estar intimamente relacionado aos outros elementos da \textbf{trilha}, em que todos devem desenhar -se mutuamente, considerando a relação de reciprocidade entre o ensino e o ato de avaliar.
Além disso, a avaliação poderá ser realizada por todas as áreas do conhecimento, de forma a analisar o desenvolvimento das habilidades propostas por cada área.
Sendo assim, a avaliação é vista como um elemento mediador, que permite conhecer as habilidades desenvolvidas pelos(as) estudantes, bem como analisar e replanejar as ações pedagógicas dos docentes envolvidos no percurso formativo.
Se o \textbf{foco}, nesta \textbf{trilha}, é estimular os(as) estudantes a pensar, a estabelecer relações, a pesquisar, a serem propositivos e criativos, os instrumentos de avaliação devem possibilitar a demonstração dessas \textbf{capacidades}.
Para isso, é importante atentar a como os(as) estudantes \textbf{sistematizam} os conhecimentos aprofundados ao longo das unidades curriculares, especialmente os que estão atrelados ao desenvolvimento gradual do projeto de pesquisa.
Avaliar esse tipo de aprendizagem, que se relaciona com a \textbf{capacidade} de ação, não é \textbf{simples} ; alguns instrumentos, porém, podem \textbf{auxiliar} nesse processo, como : produção de textos em diferentes gêneros do discurso, a exemplo de artigos científicos, artigos de opinião, relatórios, \textbf{seminários}, \textbf{debates}, \textbf{entrevistas}, projetos, portfólios ; da mesma forma, instrumentos que possibilitem interpretações de texto, resolução de exercícios, dinâmicas, experimentos, registros no caderno, \textbf{provas} escritas e orais, produção de textos em mídias \textbf{digitais} e produção de mapas mentais que interligam diferentes áreas do conhecimento.
Além disso, é muito importante que a autoavaliação seja considerada ao longo do percurso e que se avalie o impacto do desenvolvimento dos projetos de intervenção na comunidade.
Ainda, ressalta -se, para que o processo \textbf{avaliativo} seja coerente com o percurso desenhado para esta \textbf{trilha}, ser fundamental que o objeto de avaliação não se \textbf{restrinja} a conteúdos \textbf{conceituais}, mas envolva conteúdos procedimentais, contemplando atividades e situações que possibilitem a avaliação de como cada \textbf{aluno} relaciona \textbf{teoria} e prática, por meio de trabalhos individuais e em equipe.

% matched lemmas: abordagem, aluno, análise, aplicativo, aprofundamento, argumento, atual, aula, autor, auxiliar, avaliativo, caminho, capacidade, categoria, central, compreensão, conceitual, conduta, constitutivo, contemporâneo, crítico, debate, demandar, diagnóstico, digital, direcionamento, divulgar, dizer, econômico, entrevista, erro, esperar, focar, foco, fragilidade, fundamento, identitária, implicar, induzir, influência, inglês, inter-relacionar, introdução, juventude, maneira, metodológico, mobilizar, multissemiótico, ordenado, persistência, pluralidade, posicionamento, potencial, problema, processar, prova, quadro, qualitativo, relevante, remeter, restringir, resultante, seminário, sequência, simples, sistematizar, sistematização, socioambiental, sugerir, tanto, tecnologia, teoria, trilha, virtual, vista
\end{textsample}
